\section{CLS 拡散界面における粘性補間の厳密な導出}
\label{app:mu_interp_derivation}
% ═══════════════════════════════════════════════════════════════════════════════

CLS 法では $\mu(\psi) = \mu_l + (\mu_g - \mu_l)\psi$ が $\psi$ の線形関数である．
フェイス $f$（位置 $x_f$）を挟むセル $L$（中心 $x_L$）と $R$（中心 $x_R$）の間で，
粘性フラックス係数の FVM 体積平均を評価する．

CLS の平衡プロファイル $\psi_{\mathrm{eq}}(\phi) = H_\varepsilon(\phi)$ は
界面近傍で滑らかに変化するが，セル幅 $h$ が $\varepsilon$ 程度のとき
$\psi$ はセル内で（近似的に）線形と見なせる（Eikonal 条件 $|\bnabla\phi|=1$ の下で
$\psi$ は $x$ に対して線形；条件が緩い場合は近似となる）．線形 $\psi$ の下では：
\[
  \bar{\mu}_f
  \equiv \frac{1}{\Delta x}\int_{x_L}^{x_R} \mu(\psi(x))\,dx
  = \mu_l + (\mu_g-\mu_l)\,\frac{1}{\Delta x}\int_{x_L}^{x_R}\psi(x)\,dx
  = \mu_l + (\mu_g-\mu_l)\,\frac{\psi_L+\psi_R}{2}
  = \frac{\mu_L+\mu_R}{2}
\]
ここで最後の等号は $\mu_k = \mu(\psi_k)$ を用いた（$k=L,R$）．

したがって，CLS 拡散界面において $\mu$ のフェイス値は\textbf{算術平均}
$\mu_f = (\mu_L+\mu_R)/2$ が FVM 体積平均として厳密に導かれる．

これに対し $\rho^{-1}$（PPE 演算子の係数 $a = 1/\rho$）に対してはフラックス連続条件（直列抵抗則）から
調和平均 $a_f = 2/(\rho_L+\rho_R)$ を用いる（付録~\ref{app:fvm_face_coeff} 参照）．\qed

\noindent\textbf{実装との対応：} 本稿の標準実装（\texttt{ns\_terms/viscous.py}）では，粘性項を FVM フェイス積分ではなく CCD 積則微分 $\nabla\cdot[\mu(\nabla\bm{u}+\nabla\bm{u}^\top)]$ として計算する（各格子点で $\mu$ と各速度勾配を CCD で求め合成して再度 CCD 微分をとる）．この実装は FVM 算術平均とは異なるアプローチだが，$\mu$ が界面近傍以外では滑らかな場合に等価な $\mathcal{O}(h^6)$ 精度の結果を与える．本付録の算術平均証明は代替 FVM 実装を選択する場合の数学的正当化として掲載している．

% ═══════════════════════════════════════════════════════════════════════════════
